\documentclass[10pt,a4paper]{amsart}
\usepackage[margin=19mm]{geometry}
\usepackage{amsmath,amssymb,mathtools}
\usepackage[scr=rsfs,cal=euler]{mathalfa}
\usepackage{xfrac}
\usepackage[unicode,hidelinks,bookmarks=false]{hyperref}
\newcommand{\ie}{i.e.\ }
\newcommand{\cf}{cf.\ }
\newcommand{\resp}{respectively\ }
\newcommand{\Subsection}{\subsection}
\providecommand{\nobreakdash}{\nobreak\hbox{-}}
\setlength{\emergencystretch}{2em}
\multlinegap0pt
\usepackage{amssymb}
\usepackage{mathtools}
\usepackage[shortlabels]{enumitem}
\newtheorem{proposition}{Proposition}
\title{finite free probability and $S$ transforms \\ of compact Jacobi processes}
\author{Nizar Demni ${}^{*}$, Nicolas Gilliers ${}^{*,\,\dagger}$, Tarek Hamdi${}^{\triangle}$}
\date{Source: arXiv:2511.02758v4 (CC BY 4.0)}
\begin{document}
\maketitle

\noindent\emph{Source and licence.} finite free probability and $S$ transforms \\ of compact Jacobi processes. Version arXiv:2511.02758v4 (CC BY 4.0), distributed by arXiv under the Creative Commons Attribution 4.0 International licence, http://creativecommons.org/licenses/by/4.0/, as recorded in the arXiv licence metadata for this exact version and shown on the abstract page https://arxiv.org/abs/2511.02758v4. Full licence notice: \texttt{licenses/arxiv-2511.02758-CC-BY-4.0.txt}.\par\smallskip

\noindent\textit{Editorial scope.} Counted unit: the article's proposition prop:heatequationjacobi (source0.tex lines 694--706). The article's complete original proof of it is delivered with it (source0.tex lines 707--791, one \texttt{proof} environment of 85 lines). No mathematics is written, repaired or re-derived: the statement and the proof are the article's own lines, in the article's own notation, and no retained line is substituted or re-flowed.  No further source-local context is needed: every cross-reference of every retained line resolves inside the two delivered spans.  Nothing was omitted from any retained span, so the package carries no omission record.  No retained line contains a citation, so the wrapper carries no bibliography and no bibliography entry.  No retained line contains a figure environment, an image-inclusion command or drawing code, so no image is delivered and the package declares no asset dependency.  Editorial formatting only, in addition: an AMS \texttt{amsart} preamble that restates the article's own declarations and macros used by the retained lines, namely the environments \texttt{proposition} on separate counters, together with the standard packages those lines need.  The retained environments print in this document as Proposition 1 in order of appearance, whereas the article prints the same items with the numbering of its own full document.  The complete native source of the article as distributed in the arXiv e-print for this exact version is bundled unchanged as \texttt{sources/188444/source0.tex}.
\begin{proposition}
	\label{prop:heatequationjacobi}
	Let $m\geq 1$. Assume $p \wedge q > m-1$, then $(t,x) \mapsto \chi_t^{(r,s,m)}(x)$ solves the inverse heat equation:
	\begin{equation*}
		\partial_t\chi_t^{(r,s,m)}(x) = -\left\{\mathcal{L}_x^{(r,s)} + m(r+s+m+1) \right\}
		\chi_t^{(r,s,m)}(x),
	\end{equation*}
	where
	\begin{equation*}
		\mathcal{L}_x^{(r,s)} := x(1-x)\partial_{xx}^2 + [(r+1) - (r+s+2)x]\partial_x
	\end{equation*}
	is the one-dimensional Jacobi operator.
\end{proposition}

\begin{proof}
	Let
	\begin{equation*}
		F(x,\lambda_t^1, \dots, \lambda_t^m) =
		\prod_{i=1}^m(x-\lambda_t^i), \quad x \in \mathbb{R},
	\end{equation*}
	be the characteristic polynomial of the Hermitian Jacobi process $X_t$ and recall that the eigenvalues are almost surely simple. Then, It\^o's formula entails:
	\begin{multline*}
		dF(x, \lambda_t^1, \dots, \lambda_t^m) = \textrm{Local Martingale}
		-F(x, \lambda_t^1, \dots, \lambda_t^m)\sum_{i=1}^m\frac{p-(p+q)\lambda_t^i}{x-\lambda_t^i}
		\\-F(x, \lambda_t^1, \dots, \lambda_t^m)\sum_{i=1}^m\frac{1}{x-\lambda_t^i}
		\sum_{j\neq i}\frac{\lambda_t^i(1-\lambda_t^j) + \lambda_t^j(1-\lambda_t^i)}{\lambda_t^i-\lambda_t^j} dt.
	\end{multline*}
	Note in passing that the quadratic variation terms have zero contribution since
	\begin{equation*}
		\partial_{u_iu_i}F(x, u_1, \dots, u_m) = 0,
	\end{equation*}
	and since the Brownian motions $(B_t^i, t \geq 0)_{1 \leq i \leq m}$ are independent.
	Using
	\begin{equation*}
		\partial_xF(x, u_1, \dots, u_m) = F(x, u_1, \dots, u_m)\sum_{i=1}^m \frac{1}{x-u_i},
	\end{equation*}
	it follows that
	%\begin{equation*}
	%R_t(x, \lambda_t^1, \dots, \lambda_t^m) \sum_{i=1}^m\frac{p-(p+q)\lambda_t^i}{x-%\lambda_t^i}
	%= [p+m(p+q) + (p+q)x]\partial_x R_t(x, \lambda_t^1, \dots, \lambda_t^m).
	%\end{equation*}

	\begin{align*}
		F(x, \lambda_t^1, \dots, \lambda_t^m) \sum_{i=1}^m\frac{p-(p+q)\lambda_t^i}{x-\lambda_t^i}
		 & = p\partial_x F(x,\lambda_t) - (p+q)F(x,\lambda_t)\sum_{i=1}^m \frac{\lambda^i_t}{x-\lambda_t^i} \\
		 & = p\partial_x F(x,\lambda_t) - (p+q)F(x,\lambda_t)\sum_{i=1}^m (\frac{x}{x-\lambda^i_{t}}-1)     \\
		 & = p\partial_x F(x,\lambda_t) - (p+q)(x\partial_x F(x,\lambda_t)-mF(x,\lambda))                   \\
		 & = \partial_x  F(x,\lambda_t)(p - (p+q)x)+m(p+q)F(x,\lambda_t)
	\end{align*}
	where we used the shorthand notation $(x,\lambda_t)$ to denote $(x,\lambda_t^1, \dots, \lambda_t^m)$.
	Furthermore, the term
	\begin{equation*}
		2\sum_{i=1}^m\frac{1}{x-\lambda_t^i} \sum_{j\neq i}\frac{\lambda_t^i(1-\lambda_t^j) + \lambda_t^j(1-\lambda_t^i)}{\lambda_t^i-\lambda_t^j}
	\end{equation*}
	may be symmetrised as
	\begin{align*}
		\sum_{i=1}^m\sum_{j\neq i} \frac{\lambda_t^i(1-\lambda_t^j) + \lambda_t^j(1-\lambda_t^i)}{\lambda_t^i-\lambda_t^j}\left[\frac{1}{x-\lambda_t^i} - \frac{1}{x-\lambda_t^j}\right]
		 & =\sum_{i=1}^m\sum_{j\neq i}\frac{\lambda_t^i(1-\lambda_t^j) + \lambda_t^j(1-\lambda_t^i)}{(x-\lambda_t^i)(x-\lambda_t^j)}
		\\& = 2 \sum_{i=1}^m\sum_{j\neq i}\frac{\lambda_t^i(1-\lambda_t^j)}{(x-\lambda_t^i)(x-\lambda_t^j)}.
	\end{align*}
	On the other hand, one readily computes:
	\begin{equation*}
		\partial_{xx}^2F(x, u_1, \dots, u_m) = F(x, u_1, \dots, u_m)\sum_{i=1}^m\sum_{j \neq i}\frac{1}{(x-u_i)(x-u_j)}
	\end{equation*}
	and split, for any $1 \leq i \neq j \leq m$,
	\begin{align*}
		\frac{x(1-x)}{(x-u_i)(x-u_j)} & = \frac{1-x}{x-u_j} + \frac{u_i(1-x)}{(x-u_i)(x-u_j)}
		\\& = \frac{1-x}{x-u_j} + \frac{u_i(1-u_j)}{(x-u_i)(x-u_j)} - \frac{u_i}{x-u_i}
		\\& = \frac{1-x}{x-u_j} + \frac{u_i(1-u_j)}{(x-u_i)(x-u_j)} - \frac{x}{x-u_i} + 1.
	\end{align*}
	Consequently, again with the shorthand notation $(x,u)$ for $(x,u_1, \dots, u_m)$, one gets:
	\begin{align*}
		x(1-x)\partial_{xx}^2F(x, u) & = [m(m-1) + (m-1)(1-x)\partial_x -
		(m-1)x\partial_x]F(x, u)                                          \\& + F(x, u) \sum_{i=1}^m\sum_{j \neq i} \frac{u_i(1-u_j)}{(x-u_i)(x-u_j)},
	\end{align*}
	whence we deduce that
	\begin{align*}
		F(x, \lambda)  \sum_{i=1}^m & \sum_{j\neq i}\frac{\lambda_t^i(1-\lambda_t^j)}{(x-\lambda_t^i)(x-\lambda_t^j)} = x(1-x)\partial_{xx}^2F(x, \lambda)                                                                                                                        -(m-1)[m+(1-2x)\partial_x]F(x, \lambda).
	\end{align*}
	Gathering all the terms, we get the following SDE for the characteristic polynomial of $X_t$:
	%\begin{multline*}
	%dR_t(x, \lambda_t^1, \dots, \lambda_t^m) = \textrm{Local Martingale} -
	%\left\{2[p+m(p+q) + (p+q)x]\partial_x R_t(x, \lambda_t^1, \dots, \lambda_t^m) \right.
	%\\ -2 x(1-x)\partial_{xx}^2
	%R_t(x, \lambda_1, \dots, \lambda_m) \\  
	%\left.+2[m(m-1) + (m-1)(1-x)\partial_x + (m-1)x\partial_x]R_t(x, \lambda_1, \dots, \lambda_m)\right\} dt.
	%\end{multline*} 

	\begin{multline}\label{SDE1}
		d[F(x,\lambda_t)]= \textrm{Local Martingale} +\\
		-x(1-x)\partial_{xx}^2F(x,\lambda_t) - (r+1-(s+r+2)x)\partial_x F(x,\lambda_t)
		- m(d-m+1)F(x,\lambda_t) %(d-(s+r))\frac{(d+s+r+2)}{2}. 
	\end{multline}
	Finally, the bracket of the local martingale part is given by
	\begin{equation*}
		2\sum_{i=1}^m \lambda_t^i(1-\lambda_t^i)\prod_{j \neq i}(x-\lambda_j),
	\end{equation*}
	and is obviously bounded for any fixed $x$ in a bounded interval. Consequently, the local martingale part is a true martingale, and the inverse Jacobi-heat equation follows after taking the expectation of both sides of \eqref{SDE1}.
\end{proof}

\end{document}
